\section{Free-basis incidence
systems}\label{free-basis-incidence-systems}

The following argument reuses the complete free-basis proof from the
earlier phase197 manuscript; its reduced-word and tree arguments are
included here in full. Classical reduced-word normal forms and the
universal property of free groups remain explicit inputs. No historical
priority is asserted. The object set may now be infinite, but the type
set remains finite.

Let \(I\) be a finite set of cardinality \(r\ge2\), let
\(F=F(s_i:i\in I)\) be the free group on this basis, and define \[
F_i=\langle s_j:j\ne i\rangle,
\qquad F_J=\bigcap_{j\in J}F_j.
\] The objects of type \(i\) are the right cosets \(F_i g\). Objects of
distinct types are incident when the cosets intersect. Denote this
incidence system by \(\Gamma_I\). Use the left permutation convention
for the natural action: \[
R_g(F_i h)=F_i h g^{-1}.
\tag{F1}
\] Thus \(g\mapsto R_g\) is a homomorphism.

\section{Free-basis cosets and the Helly
condition}\label{free-basis-cosets-and-the-helly-condition}

A standard free factor generated by a subset of the basis consists
exactly of the elements whose reduced words use only that subset. In
particular \[
F_J=\langle s_k:k\in I\setminus J\rangle.
\tag{F2}
\] This includes \(F_\varnothing=F\) and \(F_I=1\).

For a reduced word \(w\) and \(i\in I\), let \(P_i(w)\) be the prefix
ending at its last occurrence of \(s_i\) or \(s_i^{-1}\); set \(P_i(w)\)
to the empty word if there is no such occurrence.

\textbf{Lemma R2.Prefix.} For \(i\notin J\), \[
w\in F_JF_i
\quad\Longleftrightarrow\quad
P_i(w)\text{ uses no letter indexed by }J.
\tag{F3}
\]

\textbf{Proof.} If \(w=ab\) with \(a\in F_J\) and \(b\in F_i\), reduce
the concatenation of their reduced words. Cancellation takes place
across their common boundary. The remaining word is a prefix of the
reduced word of \(a\) followed by a suffix of that of \(b\), with no
further cancellation. The second part has no occurrence indexed by
\(i\), so the prefix ending at the last such occurrence lies entirely in
the first part. It avoids \(J\). This also covers complete cancellation
of the first part, when the prefix is empty. Conversely, if the
specified prefix avoids \(J\), split \(w\) immediately after its last
\(i\)-occurrence. The prefix belongs to \(F_J\), and the suffix has no
\(i\)-occurrence and belongs to \(F_i\). \(\square\)

\textbf{Corollary R2.Product.} For every \(J\subseteq I\) and
\(i\notin J\), \[
F_JF_i=\bigcap_{j\in J}F_jF_i.
\tag{F4}
\]

\textbf{Proof.} By (F3), membership in the right side says that
\(P_i(w)\) avoids each \(j\in J\), which says exactly that it avoids
\(J\). The empty intersection causes no exception. \(\square\)

We prove flag-transitivity directly, so no outside coset criterion is a
premise here. Consider the Cayley tree with edges joining \(g\) to
\(s_i^{\pm1}g\). Every right coset of a standard free factor is a convex
subtree, because the unique reduced path between its vertices uses only
its allowed generators. A finite family of pairwise intersecting
subtrees of a tree has a common vertex. For example, root the tree and
select, among the nearest-to-root vertices of the subtrees, one at
maximal depth. For any other subtree, take a point in its intersection
with the selected subtree. Both nearest vertices lie on the root path to
that point; the deeper one lies in both subtrees. It therefore lies in
every member of the family. This proves that every flag of \(\Gamma_I\)
consists of cosets with a common representative. Right translation sends
that representative to the identity and the flag to the standard flag of
its type. In particular every flag extends to a chamber.

\textbf{Theorem R2.Geometry.} The incidence system \(\Gamma_I\) is a
flag-transitive, residually connected geometry. The natural group
\(R(F)\) acts regularly on chambers.

\textbf{Proof.} The common-representative argument proves geometry and
flag-transitivity. The stabilizer of the standard chamber is
\(\bigcap_i F_i=1\) by (F2). Thus the chamber action is regular and the
natural action is faithful.

Translate an arbitrary flag to the standard flag of some type \(J\). Its
chambers are uniquely indexed by \(F_J\), and its remaining objects
identify with the cosets of \[
F_J\cap F_i=F_{J\cup\{i\}},\qquad i\notin J,
\] inside \(F_J\). To verify this identification, intersect each
incident coset with \(F_J\); it is nonempty by the common-representative
property and is a coset of the displayed intersection subgroup. Distinct
right cosets remain distinct after nonempty intersection.

The residue is consequently the same free-basis construction on the
remaining basis \(\{s_i:i\notin J\}\). If at least two types remain, its
chamber graph is connected: chambers labeled \(h\) and \(s_i^m h\) share
all objects except possibly type \(i\), and the remaining basis elements
generate \(F_J\). Each chamber is a clique with at least two vertices,
and consecutive chambers have a common object. The union of this
connected family of cliques is the entire residue incidence graph. Thus
every residue of rank at least two is connected. \(\square\)

This proof also yields exactly the generation condition quoted as
Theorem 7.4 in the source: \[
F_J=\left\langle F_{J\cup\{i\}}:i\in I\setminus J\right\rangle
\quad\text{when }|I\setminus J|\ge2.
\tag{F5}
\] Each remaining basis letter belongs to one of those factors by
choosing a different remaining index.

\textbf{Proof correction.} Membership in \(F_jF_i\) does not imply that
the entire reduced word avoids \(s_j\). Its second factor is allowed to
contain that letter. Theorem 7.3's printed proof uses this implication.
Formula (F3), rather than avoidance by the whole word, supplies the
valid common argument.

\section{Palette actions of a free
basis}\label{palette-actions-of-a-free-basis}

Let \(C=C(\Gamma_I)\) and \(K=K(\Gamma_I)\) denote the full correlation
and type-preserving groups. An automorphism \(\alpha\) of \(F\)
satisfying \(\alpha(F_i)=F_{\sigma(i)}\) induces the correlation \[
F_i g\longmapsto F_{\sigma(i)}\alpha(g).
\tag{F6}
\] The formula is well-defined because equality of cosets is preserved
by \(\alpha\), and incidence is preserved because \(\alpha\) preserves
nonempty intersection.

\textbf{Theorem R2.TypeImage.} The action of the full correlation group
on the palette has image \(\operatorname{Sym}(I)\). Thus \[
C/K\cong\operatorname{Sym}(I).
\tag{F7}
\]

\textbf{Proof.} Every permutation \(\sigma\) of the free basis extends
to an automorphism of \(F\) by the universal property of the free group,
and it sends \(F_i\) to \(F_{\sigma(i)}\). Formula (F6) therefore
realizes every palette permutation. The kernel of the palette action is
precisely \(K\) by definition. \(\square\)

There is no claim here that (F7) determines \(K\) or the entire group
\(C\). The natural subgroup \(R(F)\) need not be the full
type-preserving group.

The inherited Lean file \texttt{BasisPermutation.lean} verifies the
generic basis-permutation automorphism, its faithful dependence on the
permutation, and the equation transporting the subgroup closure
generated by every complementary basis. This is the formal
group-theoretic step in the proof of (F7). The coset quotient
construction and the inference of the full type quotient are written
proofs, not part of that file's formal scope.

\section{The normalizer of the chamber-regular
group}\label{the-normalizer-of-the-chamber-regular-group}

Write \[
W_I=C_2^I\rtimes\operatorname{Sym}(I)
\] for the signed basis permutations. Each maps \(s_i\) to
\(s_{\sigma(i)}^{\epsilon_i}\) with \(\epsilon_i\in\{1,-1\}\) and acts
by (F6).

\textbf{Theorem R2.Normalizer.} The normalizer of the natural
chamber-regular group is \[
N_C(R(F))=R(F)\rtimes W_I.
\tag{F8}
\] The signed basis subgroup is precisely the stabilizer of the standard
chamber inside this normalizer.

\textbf{Proof.} If \(c\) normalizes \(R(F)\), compose \(c\) with the
unique natural translation taking its image of the standard chamber back
to that chamber. It suffices to consider a normalizing correlation \(d\)
fixing the standard chamber. Conjugation by \(d\) induces an
automorphism \(\alpha\) of \(F\), since \(R:F\to R(F)\) is an
isomorphism. Let \(\sigma\) be the palette permutation of \(d\).

The pointwise stabilizer in \(R(F)\) of the standard chamber with type
\(i\) deleted is \[
R\!\left(\bigcap_{j\ne i}F_j\right)=R(\langle s_i\rangle).
\] Since \(d\) fixes the chamber as a set and permutes its types by
\(\sigma\), it sends this stabilizer to
\(R(\langle s_{\sigma(i)}\rangle)\). Thus \(\alpha\) maps each infinite
cyclic subgroup \(\langle s_i\rangle\) onto
\(\langle s_{\sigma(i)}\rangle\). An isomorphism of infinite cyclic
groups sends a generator to a generator, so
\(\alpha(s_i)=s_{\sigma(i)}^{\pm1}\). Hence \(\alpha\in W_I\).

A normalizing correlation fixing the standard chamber is determined by
\(\alpha\). Indeed every chamber is its unique translate by an element
of \(R(F)\), so its image is forced by conjugation; every object belongs
to a chamber. Conversely every signed basis permutation acts by (F6),
fixes the standard chamber, and normalizes the natural group. The
chamber-regular subgroup has trivial intersection with that stabilizer,
and every normalizer element has the unique factorization already
described. This proves (F8). \(\square\)

In particular, the geometry gives a weak representation of the specified
pair \[
(F,\ F\rtimes W_I)
\] when the inner group is written first as in the source's pair
notation. This is a positive statement about compatible injections into
\((K,C)\); it is not a representation of the full abstract group \(F\)
in the source's special sense.

\textbf{Corollary R2.NaturalObstruction.} For \(r\ge2\), there is no
injective homomorphism \[
\operatorname{Aut}F\longrightarrow C(\Gamma_I)
\] whose restriction to \(\operatorname{Inn}F\cong F\) is the natural
embedding \(R\).

\textbf{Proof.} The inner subgroup is normal, so the image of such a
homomorphism would lie in \(N_C(R(F))\). Injectivity, and equality of
the image of the inner subgroup with \(R(F)\), would induce an injection
\[
\operatorname{Out}F\hookrightarrow N_C(R(F))/R(F)\cong W_I.
\] The right side is finite. The left side is infinite: for distinct
basis elements \(s_1,s_2\), the automorphisms \(s_1\mapsto s_1s_2^m\),
fixing the other basis elements, induce distinct transformations of the
abelianization \(\mathbb Z^r\), whereas inner automorphisms act
trivially there. This is a contradiction. \(\square\)

The corollary fixes the natural embedding used in the construction. It
does not rule out every conceivable weak embedding with a different
image of \(\operatorname{Inn}F\). The source's finite-rank definition
separately rules out an \emph{exact} representation of a nonabelian free
group: its outer quotient would have to embed into the finite palette
permutation group. That observation is already made in {[}L126, §7{]}.

\section{Weak representations from panel permutations}
The natural-normalizer obstruction fixes the prescribed chamber-regular
inner action. It does not obstruct other compatible embeddings in the
source's definition of a weak representation. The following construction
settles that distinction without changing the incidence geometry.

\begin{theorem}[An unrestricted panel-permutation subgroup]
\label{UR.PanelSymmetric}
For a finite free basis of size $r\ge2$, the type-preserving group of
$\Gamma_I$ contains $\Sym(\mathbb Z\setminus\{0\})$. Consequently
it contains every countable group.
\end{theorem}
\begin{proof}
Fix $i\in I$. Write elements of $F$ in their unique reduced syllable
form $s_{j_1}^{n_1}\cdots s_{j_k}^{n_k}$, with nonzero integer
exponents and consecutive indices different. Extend any
$\pi\in\Sym(\mathbb Z\setminus\{0\})$ by $\pi(0)=0$. Define
$\Phi_\pi:F\to F$ as follows: if the final syllable is $s_i^n$,
replace it by $s_i^{\pi(n)}$; otherwise leave the word unchanged.
The empty word is fixed. Nonzero exponents remain nonzero, so reduced
syllable forms stay reduced. It follows directly that
$\Phi_\pi\Phi_\rho=\Phi_{\pi\rho}$ and
$\Phi_\pi^{-1}=\Phi_{\pi^{-1}}$.

For $j\in I$, call $g,h$ $j$-adjacent when
$h=s_j^m g$ for some integer $m$; these are precisely the chambers in
one $j$-panel. The map $\Phi_\pi$ preserves this relation in both
directions. On the panel $\langle s_i\rangle$ this is immediate: it
is an arbitrary permutation of that panel fixing its identity chamber.
Outside that case, multiplication on the left changes or cancels only
the first syllable. It leaves the final $i$-syllable unchanged, if one
is present; even when a two-syllable word loses its first syllable,
both orders of the two operations give $s_i^{\pi(n)}$. If no final
$i$-syllable is present the words remain fixed, except for the already
handled panel at the identity. Thus adjacent images are again in the
same $j$-panel. Applying the inverse proves the converse.

The chambers in an object $F_jg$ are exactly those reachable from $g$
using panels of labels other than $j$, because $F_j$ is generated by
the corresponding cyclic factors. Therefore $\Phi_\pi$ permutes the
right cosets of every $F_j$. Define its object action by
$F_jg\mapsto F_j\Phi_\pi(g)$. The preceding observation proves
well-definedness and bijectivity. Two objects are incident exactly when
their sets of chambers intersect, so this action preserves incidence
and types. Distinct permutations $\pi$ give distinct object actions: the type
$i$ cosets $F_is_i^n$ are distinct, since
$F_i\cap\langle s_i\rangle=1$, and their images detect $\pi(n)$.
This proves the symmetric-group embedding.

Finally a countable group embeds faithfully in the symmetric group of
its underlying set by its regular action. Identify that set with a
subset of $\mathbb Z\setminus\{0\}$ and fix the complement to obtain
the asserted embedding.
\end{proof}

\begin{corollary}[Weak representation of a finite-rank free group]
\label{UR.WeakFree}
The geometry $\Gamma_I$ is a weak incidence representation of the
abstract free group $F$ in the source's precise sense. Both compatible
images may be taken inside its type-preserving group.
\end{corollary}
\begin{proof}
An automorphism of $F$ is determined by the images of its finitely many
basis elements. The set $F^r$ is countable, so $\Aut(F)$ is countable.
Embed $\Aut(F)$ in $K(\Gamma_I)$ using
Theorem~\ref{UR.PanelSymmetric}, and restrict this injection to
$\operatorname{Inn}(F)$. Composing the first map with
$K(\Gamma_I)\hookrightarrow C(\Gamma_I)$ gives the required commuting
pair of injections.
\end{proof}
This is an existence construction with a different inner image. It
neither extends the natural chamber-regular action nor contradicts
Theorem R2.Normalizer and its natural-embedding obstruction. No quotient
identification is inferred from these weak injections.
